\subsection{The ideal-class map and the principal divisor of a secondary unit lift}

\begin{statement}{A.208}{Proved}
{An exact support-class isomorphism, principalization, and next norm matrix}
Use the embedded fields, units, and automorphisms of A.203 and A.206.
Write \(U_L=\mathcal O_L^\times\),
\(\widehat U_L=\widehat{\mathcal O_L^\times}_3\),
\(A=\operatorname{Cl}(M)\{3\}\), and
\(A_2=\operatorname{Cl}(E_2)\{3\}\).
The actual unit-support quotient has the canonical isomorphism
\[
 \Theta_h:U_{E_2}/N_hU_M
       \xrightarrow{\ \sim\ }A_2^-,\qquad
 \Theta_h(a)=N_h([D]_3),\quad
 N_h(w)=a,\quad \operatorname{div}(w)=(h-1)D.
\tag{A.208.1}
\]
Here \(D\) is an actual fractional ideal divisor of \(M\),
the subscript \(3\) denotes its three-primary class, and
\(A_2^-\) is the minus eigenspace for complex conjugation.
The same map identifies the completed unit quotient with \(A_2^-\).
Consequently the A.203 map gives the specified onto morphism
\[
 \Xi=\Theta_h\circ\eta:\mathcal R_1\longrightarrow A_2^-,
 \qquad
 \ker\Xi=\ker\eta.
\tag{A.208.2}
\]

Put \(\Omega=K_3(z)\), where \(z=\zeta_9\), and retain
the actual pure roots
\[
 \rho_1=d^8\alpha,\quad \rho_2=d^2t,\quad
 \rho_3^3=\rho_2,\quad
 \gamma(\rho_3)=z\rho_3,\quad \gamma(z)=z,
 \qquad s=\gamma^3.
\tag{A.208.3}
\]
Then \(\Omega^{\langle\gamma\rangle}=F'\),
\(\Omega^{\langle s\rangle}=M\), and \(\gamma|_M=\sigma\).
For each actual anti-unit \(a\in U_{E_2}\), set
\(v=\sigma^2(a)/\sigma(a)\).  Every actual pure-unit preimage
\(b\in U_{K_3}\) with \(N_{K_3/K_2}b=v\) has a factorization
\[
 b=\frac{\gamma(u)}u\frac{\sigma(x)}x,
 \qquad u\in U_\Omega,\quad x\in M^\times,\quad
 J=\operatorname{div}_M(x),\quad \sigma J=J.
\tag{A.208.4}
\]
Such actual preimages exist by the one-step unit norm theorem A.184.
The construction below gives \(3J=\operatorname{div}_M(f)\)
for an actual \(f\in F'^\times\), and proves the full ordinary
ideal class \([J]\) is zero before choosing \(x\).
At each \(q=11,59\), let \(\mathfrak R_q\) be the prime of
\(M\) whose residue map has \(\alpha=\alpha_{0,q}\zeta\),
where \(\alpha_{0,q}\in\Q_q\) is the unique cube root of three
and \(\zeta=z^3\).  With the precise A.184 normalization,
\[
 \psi_2(v)_q=2v_{\mathfrak R_q}(x)\pmod3.
\tag{A.208.5}
\]
This is the character in the pure residue field \(\F_{q^2}\).

The complete pro-three unit group of \(M\), with all 51 A.199
generators, has a next norm-character image of dimension three:
\[
 [U_M:U_M\cap N_s\Omega^\times]=3^3,
 \qquad
 \widehat U_M/
   (\widehat U_M\cap N_s\widehat{\Omega^\times}_3)
       \simeq(\Z/3)^3.
\tag{A.208.6}
\]
The six literal tame rows specifying this quotient are displayed
below.  In particular
\[
       |\operatorname{Cl}(\Omega)\{3\}^{\langle s\rangle}|
                  =3^7.
\tag{A.208.7}
\]
Equations (A.208.1) and (A.208.5) retain their respective ideal
constructions.  Their comparison is the remaining explicitly
described alignment problem at the end of the proof.
\end{statement}

\begin{proof}
\emph{The unit-support class and all its choices.}
A.203 proves that every actual \(a\in U_{E_2}\) is an actual
field norm from \(M\).  Choose \(w\) of that norm.  Its
residue-weighted ideal norm is zero.  On a split orbit of the
cyclic group \(\langle h\rangle\), the coefficients of
\(\operatorname{div}(w)\) sum to zero, and this is exactly the
image of \(h-1\) on the integral orbit lattice.  On a nonsplit
orbit its single coefficient is zero because its positive
residue degree multiplies it to zero.  Thus there is an actual
fractional ideal divisor \(D\) with
\(\operatorname{div}(w)=(h-1)D\).  Its ideal class is
\(h\)-invariant.

If \(w\) is replaced by another preimage, their ratio is
\(h(y)/y\) by multiplicative Hilbert--90.  Therefore \(D\)
changes by a principal divisor and an \(h\)-invariant divisor.
The latter is an extended divisor plus an integral multiple
of the unique ramified prime \(\mathfrak P_3\) of \(M/E_2\).
A.206 gives
\[
 j_hA_2=0,\quad
 \ker(N_h:A\to A_2)=\langle c_3\rangle,
 \quad |A^h|=27,
 \qquad c_3=[\mathfrak P_3]_3.
\tag{A.208.8}
\]
Hence the primary classes of all those invariant divisors
belong to \(\langle c_3\rangle\), and their ideal norms
vanish in \(A_2\).  This proves independence in (A.208.1),
with the sign fixed by the displayed \(h-1\) convention.
Multiplication of norm preimages proves additivity of the map.

The intermediate map is precisely
\[
 U_{E_2}/N_hU_M\longrightarrow A^h/\langle c_3\rangle,
                   \qquad a\longmapsto[D]_3.
\tag{A.208.9}
\]
To prove surjectivity, represent any \(h\)-fixed primary
class by a fractional ideal divisor \(D\).  The divisor
\((h-1)D\) is principal; a generator \(w\) has unit norm,
so its unit norm maps to the chosen class.  For injectivity,
suppose \([D]_3\in\langle c_3\rangle\).  Choose an
\(h\)-invariant divisor \(I\), supported on
\(\mathfrak P_3\), with \([D-I]_3=0\).  There is an
integer \(e\) prime to three and \(y\in M^\times\) with
\(e(D-I)=\operatorname{div}(y)\).  The literal element
\[
        w^e\big/\bigl(h(y)/y\bigr)\in U_M
                  \quad\hbox{has norm }a^e.
\tag{A.208.10}
\]
Choose integers \(k,l\) with \(ke+3l=1\), and multiply
its \(k\)-th power by the scalar unit \(a^l\).  The
result is an actual unit of norm \(a\).  Thus (A.208.9)
is an isomorphism before completion.

Ideal norm identifies its target with its image in \(A_2\),
by (A.208.8).  Complex conjugation \(\tau\) commutes with
\(h\), so every map just constructed is \(\tau\)-equivariant.
The A.203 quotient is generated by the anti-units \(r_3,r_4\),
and has order nine.  Its image is therefore the entire
two-dimensional \(A_2^-\) of A.173.  This proves
(A.208.1).  The actual quotient is killed by three, so passage
to pro-three completions gives the same finite quotient and
the same map.  Composition with the literal difference map
and quotient of A.203 proves (A.208.2).

\emph{The cyclic-nine field, with its actual generators.}
The relations \(\rho_2=d^2t\) and
\(\rho_2^3=d^8\alpha\) prove all the fixed fields in
(A.208.3).  The automorphism denoted \(g\) in A.199 acts on
\(M\) by \(t\mapsto zt\) and moves \(\alpha\) to
\(\zeta\alpha\).  In contrast, restriction of the present
\(\gamma\) fixes \(\alpha\) and is \(\sigma\).
These assertions follow from the actual powers of the two
displayed root coordinates.  Extend \(h\) to \(\Omega\)
by fixing \(\rho_3\) and sending \(z\) to \(z^4\).
The literal relations are
\[
 h\gamma h^{-1}=\gamma^4,
 \qquad \tau\gamma\tau=\gamma^{-1},
 \qquad \tau h=h\tau,
 \qquad hs=sh.
\tag{A.208.11}
\]
The last equality uses \(s^4=s\).  At every prime above
\(3,11,59\), both \(M/F'\) and \(\Omega/M\) are
totally ramified cubic extensions.  At three this is the
full local root-degree calculation of A.140 and A.162.
At each changed prime the valuations of \(\rho_1\) and
\(\rho_2\) in their respective bases are eight, prime to
three.  Away these primes, the root polynomial and its
derivative show that these extensions are unramified.

\emph{Principalizing the entire Hilbert--90 divisor.}
For an actual anti-unit \(a\), the element
\(v=\sigma^2(a)/\sigma(a)\) is fixed by \(\tau\)
and has \(\sigma\)-norm one.  Thus \(v\in U_{K_2}\)
and \(N_{K_2/K_1}v=1\).  A.184 first kills its two tame
characters and then supplies an actual \(b\in U_{K_3}\)
of norm \(v\).  The norm \(N_\gamma b=N_\sigma v\)
is one.  Hilbert--90 gives \(w\in\Omega^\times\) with
\(b=\gamma(w)/w\).  Write \(m=N_sw\).  The equality
\(\sigma(m)/m=v\) is equivalent to
\[
                    m=f\sigma(a),\qquad f\in F'^\times.
\tag{A.208.12}
\]
The divisor \(D=\operatorname{div}_\Omega(w)\) is
\(\gamma\)-invariant, since \(b\) is a global unit.
At any prime ramified in \(\Omega/M\), its norm coefficient
is its own coefficient, since the residue degree is one.
By (A.208.12) this equals the coefficient of
\(\operatorname{div}_M(f)\), divisible by three because
\(M/F'\) is totally ramified there.  At the unramified
primes an invariant divisor descends without a divisibility
condition.  It follows, at each actual prime, that
\[
 D=j_sJ,\qquad \sigma J=J,\qquad
                   3J=\operatorname{div}_M(f)
\tag{A.208.13}
\]
for an actual fractional ideal divisor \(J\) of \(M\).
The last equality follows by the literal norm identity
\(N_sj_s=3\) on divisors.

The primary classes of \(\sigma\)-invariant divisors are
generated by classes extended from \(F'\) and by the primes
ramified in \(M/F'\).  The former have zero primary class
because \(\operatorname{Cl}(F')\{3\}=0\).  At the changed
primes, \(M/E_2\) is unramified inert: its residue degree is
three because \(q^2\) has order three modulo nine.  Their
classes are therefore literal extensions of lower classes,
whose primary parts vanish by A.206.  Consequently
\([J]_3\in\langle c_3\rangle\).  This uses no assertion
about inertness of the degree-six primes above two.

There is also an actual conjugation identity.  Since
\(\tau(b)=b\), the element
\(a_0=\gamma(w)\tau(w)\) is \(\gamma\)-fixed.  Indeed
\(\tau\gamma=\gamma^{-1}\tau\), applied to
\(\gamma(w)/w=b\), makes the two factors of
\(\gamma(a_0)/a_0\) inverse.  Hence \(a_0\in F'^\times\),
and
\[
 a_0^3=N_s(a_0)=\sigma(m)\tau(m)=f\tau(f),
 \qquad J+\tau J=\operatorname{div}_M(a_0).
\tag{A.208.14}
\]
The middle cancellation uses exactly \(\tau a=a^{-1}\)
and \(\tau\sigma=\sigma^{-1}\tau\).  The divisor identity
follows from (A.208.13), since the divisor group is torsion-free.
The unique prime \(\mathfrak P_3\) is \(\tau\)-fixed.
Thus \([J]_3\) is both in its plus line and in the minus
eigenspace, so it is zero.  Equation (A.208.13) already kills
the full class \([J]\) by three.  Therefore \([J]=0\)
in the ordinary class group, and an actual \(x\) with
\(\operatorname{div}(x)=J\) exists.  Put \(u=w/x\).
Its divisor is zero, and direct substitution proves
(A.208.4), including \(\sigma(x)/x\in U_M\).

\emph{The exact secondary residues.}
At the specified \(\mathfrak R_q\), the element
\(\pi_M=t^2/q\) is a uniformizer, since
\(v(t)=2\) and \(v(q)=3\).  It satisfies
\(\sigma(\pi_M)/\pi_M=\zeta^2\).
The extension \(\Omega/F'\) is totally ramified at this
prime.  Its generator \(\gamma\) therefore acts trivially
on the residue of the unit \(u\).  Equation (A.208.4) gives
\[
        \overline b=\zeta^{2k_q},
              \qquad k_q=v_{\mathfrak R_q}(x).
\tag{A.208.15}
\]
Here the residue \(\overline b\) lies in the pure residue
field \(\F_{q^2}\), embedded in the \(\F_{q^6}\) residue
field of \(M\).  Put \(E_q=(q^2-1)/3\).
Since \(\alpha_{0,q}\in\F_q^\times\) and
\(q-1\) divides \(E_q\), the selected branch has
\(\lambda_q(\alpha)=\zeta^{E_q}\).  The two values
\(E_{11},E_{59}\) are respectively \(1,2\) modulo three.
Consequently
\(\lambda_q(b)=(\lambda_q(\alpha))^{2k_q}\), proving
(A.208.5) with the A.184 character exactly.
Changing \(w\) by an \(F'^\times\) element changes \(k_q\)
by a multiple of three; changing \(x\) by a unit changes no
valuation or ramified residue.  Independence from the selected
pure-unit preimage is precisely the actual norm-kernel theorem
in A.184.  The continuous completed map is its unique extension
from the finitely generated unit lattice.

\emph{The complete next tame norm matrix.}
For \(q=11,59\) the residue field of each prime of \(M\)
is \(\F_{q^6}=\F_q[z]/(z^6+z^3+1)\), since the order
of \(q\) modulo nine is six.  Label its three primes by
\(\alpha=\alpha_{0,q}z^{3i}\), \(i=0,1,2\).  The
\(\Omega/M\) unit norm character is
\[
 \Lambda_{q,i}(w)=
   \log_\zeta\bigl(\overline w^{(q^6-1)/3}\bigr),
       \qquad w\in U_M.
\tag{A.208.16}
\]
Use the actual A.199 generator order and write \(g(t)=zt\),
\(g(z)=z\), solely for the orbits in that order.  All 24
columns from the three conjugate \(E_2\)-unit groups are zero.
They come from residue fields \(\F_{q^2}\), whose inclusion
in \(\F_{q^6}\) makes this cubic character zero.
The nine \(g^j(r)\) columns have period three, and all nine
\(g^j(s)\) columns are equal; these latter equalities are
verified on their fixed literal coefficient vectors.
Thus the complete matrix is specified by its following 13
columns, in the order
\[
                  e_1,\ldots,e_9,\ r,g(r),g^2(r),\ s:
\tag{A.208.17}
\]
\[
\begin{array}{c|rrrrrrrrr|rrr|r}
 (q,i)&1&2&3&4&5&6&7&8&9&r&gr&g^2r&s\\\hline
 (11,0)&1&0&1&1&1&2&2&0&2&0&2&2&2\\
 (11,1)&1&0&1&1&2&1&1&1&2&2&2&0&2\\
 (11,2)&1&0&1&1&1&2&1&0&2&2&0&2&2\\\hline
 (59,0)&2&0&2&2&1&2&0&2&1&2&0&0&1\\
 (59,1)&2&0&2&2&0&0&1&1&1&0&0&2&1\\
 (59,2)&2&0&2&2&1&2&1&2&1&0&2&0&1
\end{array}
\tag{A.208.18}
\]
If \(R_{11},R_{59}\) are its three-row blocks, then
\[
 R_{59}=\begin{pmatrix}1&2&2\\2&1&2\\2&2&1\end{pmatrix}R_{11},
 \qquad
 \det\begin{pmatrix}1&1&2\\1&2&1\\1&1&1\end{pmatrix}=-1.
\tag{A.208.19}
\]
The minor uses the \(e_1,e_5,e_7\) columns.  These equalities
prove that the full matrix has rank exactly three.

For reproducibility, the calculation uses rational coefficients
in the literal presentation (A.199.3), without a numerical local
embedding.  Write a unit as \(C_0+tC_1+t^2C_2\) by substituting
\(t^3=d^2\alpha\).  The three summands have valuations in
the three distinct classes \(0,2,1\) modulo three.  Integrality
therefore makes its residue precisely the residue of \(C_0\).
For a common rational denominator of \(q\)-valuation \(e\),
lift the unique scalar cube root \(\alpha_{0,q}\) modulo
\(q^{e+1}\), evaluate \(C_0\)'s numerator in
\((\Z/q^{e+1})[z]/(z^6+z^3+1)\), divide its coefficients
by \(q^e\), and reduce modulo \(q\).  The certificate checks
every divisibility and all 306 resulting unit values.  The
largest scalar precision is six.  It also checks the second
unit's cubic equation on every branch and every orbit element.
The literal finite unit formula for \(\epsilon\) is evaluated
in the same field.  A.199 supplies completeness of these 51
source units, so the rank computation describes the entire
unit quotient, rather than just their subgroup.

The six changed primes are totally ramified tame cubic places
of \(\Omega/M\); their unit norm conditions are exactly
(A.208.16).  There is one additional ramified place, above
three.  A global unit has zero local norm symbol at every
other finite place, since the extension is unramified there;
all infinite places are complex.  Global reciprocity makes its
symbol at the remaining three-place zero when the six displayed
symbols vanish.  The cyclic Hasse norm theorem
\cite[Chapter VIII, Theorem 3.1, printed p.~234]{MilneCFT}
then gives an actual field norm.  This proves the first quotient
in (A.208.6).  The completed assertion follows through the same
finite character map: choose actual norm preimages of generators
of its actual unit kernel for one inclusion, and localize any
completed norm for the other.  The finite quotient is unchanged.
Finally the ordinary ambiguous class formula already proved and
used in A.175 gives
\[
 |\operatorname{Cl}(\Omega)\{3\}^{s}|
    =|A|\,3^{7-1}/3^3=3^4\,3^6/3^3=3^7,
\tag{A.208.20}
\]
with the seven actual ramified primes just listed.

\emph{The retained alignment equation.}
For the three certified anti-units \(a=r_2,r_3,r_4\),
their literal coefficient vectors also give \(N_\sigma a=1\).
For their products, the condition \(h(b)=b\), applied before
principalization, gives the actual element
\[
 K=\frac{h(w)\gamma h(w)\gamma^2h(w)\gamma^3h(w)}{w}
       \in F'^\times,\qquad
 K^3=\frac{h(f)^4}{f}.
\tag{A.208.21}
\]
The first equality follows by taking \(\gamma(K)/K\) and
using (A.208.11); the second follows by applying \(N_s\)
to (A.208.12) and using \(N_\sigma a=1\).
Putting \(r_0=K/h(f)\) gives the literal identities
\[
 r_0^3=h(f)/f,\qquad
 (h-1)J=\operatorname{div}_M(r_0),\qquad
                   h(x)/(r_0x)\in U_M.
\tag{A.208.22}
\]
Thus the class map in (A.208.1) and the valuation formula
(A.208.5) come with this remaining actual unit correction.
The asserted results retain that correction without assigning
it an unevaluated value in either quotient.

The executed certificate
\path{code/two_inert_ninth_root_next_norm_characters.py}
retains all exact unit residues, the expanded six-by-51 matrix,
the compact matrix, its row relation and nonzero minor, the
literal three anti-unit norms, and the character-normalization
checks.  Its output and rational input conventions are retained.
\end{proof}
